\noindent\textbf{Motivation.} Personalized neoantigen vaccines depend on
selecting, from many candidate mutations, the few that will provoke a strong
T-cell response. A large family of immunogenicity predictors now decides whether
a peptide is recognized at all, but \emph{how strong} the response will
be---its magnitude---is essentially unaddressed, even though graded experimental
readouts such as ELISpot spot-forming counts make this quantity measurable. It is
not known how well existing tools track response magnitude when their scores are
repurposed as continuous predictors, nor whether the apparent differences between
them are real.

\noindent\textbf{Results.} We assemble a unified benchmark that places eight
immunogenicity tools and one presentation proxy on a common footing through a
seven-step harmonization protocol, and scores them against a continuous ELISpot
magnitude readout as well as a binary recognition label on a common patient
cohort (DS2, $101$ peptides). Three findings emerge. First, the correlation with
response magnitude is uniformly weak: every tool falls below an absolute Spearman
correlation of $0.33$, and binary-discrimination differences, while larger in
point estimate, have broadly overlapping bootstrap intervals that shift with
aggregation and threshold, so no single tool is statistically distinguishable as
best at the top of the panel. Second, a per-patient evaluation
protocol---correlating within each patient before aggregating---surfaces
within-patient ranking ability in a subset of patients that the default global
metric conceals, together with large individual variation. Third, we quantify
the magnitude gap and argue, with an explicit biological ceiling, that it is an
open opportunity rather than a closed problem.

\noindent\textbf{Conclusion.} Response magnitude remains an unfilled gap with a
genuine biological upper bound; our reproducible benchmark and per-patient
protocol provide a measuring stick for the continuous-magnitude predictors that
this gap invites.
